\section{Possibilities}

\subsection{Control flow}
    One of the motivations behind typed assembly languages as presented earlier
    (TODO: REF SECTION) is control flow.
    TAL-0 introduces types to code to restrict where jumps are allowed to
    target.
    Control-flow in regards to execution time analysis is not so interested in
    the validity of jumps but on the amount of jumps, i.e. loops.

    Control-flow graph

    By using DTAL and considering a program an oriented graph where blocks are
    nodes and branches are edges to create something like a control-flow graph.

\subsection{Typing memory}
    The chosen platform is explicit in its memory hierarchy, which enables it
    to be explicit in code as well. This can be done through a type system, so
    any array or stack type is wrapped inside a memory region context.

\subsection{}
